\documentclass[12pt]{article}
\usepackage[T1]{fontenc}
\usepackage[utf8]{inputenc}
\usepackage{lmodern}
\usepackage{textcomp}
\usepackage{enumitem}
\usepackage{geometry}
\geometry{margin=1in}
\begin{document}
\begin{center}
Answer Key - Science - K-12 Level Assessment 2 \
\small (Answers are ordered exactly as in the question paper.)
\end{center}

\section*{Multiple Choice}
\begin{enumerate}[label=\arabic*.]
\item \textbf{Answer:} D
\\ \textit{Options:} A) a force that pushes straight up; B) a force that pushes straight down; C) a force that pushes the same direction that the car is moving; D) a force that pushes the opposite of the direction that the car is moving
\\ \textit{Note:} A car in motion will continue to roll due to inertia, but a force acting in the opposite direction, such as friction or a push, can counteract this motion and bring the car to a stop. This aligns with Newton's first law of motion, which states that an object will remain in its state of motion unless acted upon by an external force.
\item \textbf{Answer:} B
\\ \textit{Options:} A) The total amount of energy increases.; B) The total amount of energy is constant.; C) The energy is destroyed through motion.; D) The amount of chemical energy increases.
\\ \textit{Note:} The total amount of energy is constant because of the law of conservation of energy, which states that energy cannot be created or destroyed but can only be transformed from one form to another, such as from chemical energy to mechanical energy in the human body.
\item \textbf{Answer:} D
\\ \textit{Options:} A) a hypothesis.; B) a conclusion.; C) an assumption.; D) an observation.
\\ \textit{Note:} The student has made an observation because they are noting a specific phenomenon in nature-the correlation between moss growth and forest density-without forming a hypothesis or drawing conclusions, which are steps beyond mere observation. This aligns with the key concept in science that observations are the initial, factual data collected through direct sensory experience.
\item \textbf{Answer:} C
\\ \textit{Options:} A) fruits; B) grains; C) cheese; D) vegetables
\\ \textit{Note:} A person whose gallbladder has been removed may experience difficulty digesting high-fat foods, such as cheese, because the gallbladder's role in storing bile, which emulsifies fats, is lost, leading to potential digestive discomfort and malabsorption.
\item \textbf{Answer:} C
\\ \textit{Options:} A) clear glass and green glass.; B) paper cups and plastic cups.; C) iron nails and aluminum nails.; D) sand and salt.
\\ \textit{Note:} A strong magnet will attract iron nails because they are ferromagnetic materials, while aluminum nails are not magnetic and will not be affected by the magnet, allowing for effective separation of the two types of nails. This demonstrates the principle of magnetism and its application in sorting materials based on their magnetic properties.
\end{enumerate}

\section*{True / False}
\begin{enumerate}[label=\arabic*.]
\setcounter{enumi}{5}
\item \textbf{Answer:} False
\\ \textit{Options:} A) True; B) False
\\ \textit{Note:} Renewable resources, such as solar and wind energy, can increase the overall supply of electricity but do not inherently reduce the demand for electricity, as society continues to consume energy for various needs.
\item \textbf{Answer:} True
\\ \textit{Options:} A) True; B) False
\\ \textit{Note:} The tilt of Earth's axis at approximately 23.5 degrees causes different parts of the planet to receive varying amounts of sunlight throughout the year; during December, the Northern Hemisphere is tilted away from the Sun, resulting in shorter days and less direct sunlight, which leads to colder temperatures.
\item \textbf{Answer:} False
\\ \textit{Options:} A) True; B) False
\\ \textit{Note:} The statement is false because stems primarily serve as conduits for transporting water and nutrients from the roots to other parts of the plant, while photosynthesis occurs in the leaves, where sunlight is used to convert carbon dioxide and water into food (glucose) for the plant.
\item \textbf{Answer:} True
\\ \textit{Options:} A) True; B) False
\\ \textit{Note:} The statement is true because the Alps, Appalachians, and Himalayas were all formed as a result of tectonic plate collisions and movements that caused the Earth's crust to fold and uplift, creating mountainous regions characterized by folded rock structures.
\item \textbf{Answer:} True
\\ \textit{Options:} A) True; B) False
\\ \textit{Note:} The Big Bang theory posits that the universe began as an extremely hot and dense point, which then rapidly expanded, leading to the formation of all matter and energy in the cosmos.
\end{enumerate}

\section*{Long Form Response}
\begin{enumerate}[label=\arabic*.]
\setcounter{enumi}{10}
\item \textbf{Answer:} Photosynthesis is the process by which plants convert sunlight into energy, primarily occurring in the chloroplasts of their cells. During this process, plants take in carbon dioxide from the air and water from the soil, using sunlight to transform these materials into glucose and oxygen. This transformation is crucial not only for the plants, providing them with the energy needed for growth and development, but also for the ecosystem, as it produces oxygen essential for the survival of most living organisms and serves as the foundation of the food chain.
\\ \textit{Note:} correct because it accurately describes the process of photosynthesis, detailing how plants convert sunlight, carbon dioxide, and water into glucose and oxygen, while also emphasizing the significance of this process for plant energy needs and the broader ecological role of oxygen production.
\item \textbf{Answer:} Proper labeling in scientific experiments is crucial, especially when studying plant growth rates, because it ensures that all variables and conditions are clearly identified. This clarity allows other researchers to accurately replicate the experiment and verify results, which is essential for building reliable scientific knowledge. If labels are unclear or missing, it can lead to confusion and potentially invalidate the findings, making it difficult for others to trust or use the data.
\\ \textit{Note:} correct because it emphasizes that proper labeling ensures clarity in identifying variables and conditions, which is essential for reproducibility and verification of results, ultimately affecting the reliability of scientific data. Clear labels prevent confusion and enhance the trustworthiness of findings for other researchers.
\end{enumerate}

\vfill
\noindent\textit{End of Answer Key}
\end{document}